% from ega4-10-fr.tex
\oldpage[IV-3]{106}
\begin{remark}[10.5.6]
\label{IV.10.5.6-fr}
L'assertion (ii) de (10.5.5) est encore valable lorsque $X$ est
\emph{localement noethérien}~: en effet, si $(V_\alpha)$ est un recouvrement de
$X$ formé d'ouverts affines (noethériens), tout sous-préschéma d'un $V_\alpha$
est un sous-préschéma de $X$~; inversement, si un sous-préschéma $Z$ de $X$ est
isomorphe au spectre d'un anneau semi-local intègre de dimension $1$, il y a
un $\alpha$ tel que $V_\alpha\cap Z$ contienne un ouvert affine $U$ de $Z$ non
réduit au point générique de $Z$, et qui est donc aussi spectre d'un anneau
semi-local intègre de dimension $1$. On conclut à l'aide de (10.3.3).
\end{remark}

\begin{proposition}[10.5.7]
\label{IV.10.5.7-fr}
Soient $X$ un préschéma localement noethérien, $Y$ une partie fermée de $X$
telle que toute partie fermée non vide de $X$ rencontre $Y$. Alors le
préschéma induit sur l'ouvert $X-Y$ est un préschéma de Jacobson.
\end{proposition}

Appliquons le critère (10.5.5, (ii)), et supposons qu'il y ait un
sous-préschéma $Z$ de $X-Y$ qui soit spectre d'un anneau semi-local intègre de
dimension $1$, le point générique $z$ de $Z$ étant donc isolé dans $Z$ (ou
dans l'adhérence $\overline Z$ de $Z$ dans $X$), et $Z$ étant distinct de
$\{z\}$. Soit $y\neq z$ un point de $Z$~; comme il n'appartient pas à $Y$, il
n'est pas fermé dans $X$, et son adhérence $\overline{\{y\}}$ dans $X$
rencontre $Y$ en un point $x\neq y$ qui est donc spécialisation de $y$.
L'existence de la chaîne
$\{x\}\subset\overline{\{y\}}\subset\overline{\{z\}}$ montre alors que la
dimension de $\overline{\{z\}}$ serait $\geq2$, et il en serait de même de la
dimension de $U\cap\overline{\{z\}}$, où $U$ est un voisinage ouvert affine
(donc noethérien) de $x$ dans $X$. Or, cela contredit le fait que le point
générique $z$ de $U\cap\overline{\{z\}}$ est isolé dans
$U\cap\overline{\{z\}}$ (10.5.3).

\begin{corollary}[10.5.8]
\label{IV.10.5.8-fr}
Soit $A$ un anneau noethérien~; pour tout élément $f$ du radical $\mathfrak R$
de $A$, l'anneau $A_f$ est un anneau de Jacobson, et l'ouvert
$\operatorname{Spec}(A)-V(\mathfrak R)$ est un schéma de Jacobson.
\end{corollary}

Si $X=\operatorname{Spec}(A)$, $Y=V(\mathfrak J)$, où $\mathfrak J$ est un
idéal de $A$, dire que toute partie fermée non vide de $X$ rencontre $Y$
équivaut ($\mathbf 0_{\mathrm I}$, 2.1.3) à dire que $Y$ contient tous les
\emph{points fermés} de $X$, ou encore que $\mathfrak J$ est \emph{contenu
dans le radical} $\mathfrak R$ de $A$. Si $f\in\mathfrak R$, l'ouvert $D(f)$
ne rencontre donc pas $V(\mathfrak R)$, et est un espace de Jacobson en vertu
de (10.5.7).

\begin{corollary}[10.5.9]
\label{IV.10.5.9-fr}
Soient $A$ un anneau local noethérien, $\mathfrak m$ son idéal maximal,
$X=\operatorname{Spec}(A)$, $Y=X-\{\mathfrak m\}$~; alors $Y$ est un schéma de
Jacobson, dont les points fermés sont les idéaux premiers
$\mathfrak p\in\operatorname{Spec}(A)$ tels que $\dim(A/\mathfrak p)=1$.
\end{corollary}

La première assertion est un cas particulier de (10.5.8)~; d'autre part les
points fermés de $Y$ sont les idéaux premiers $\mathfrak p$ de $A$ qui sont
des éléments maximaux dans l'ensemble des idéaux premiers $\neq\mathfrak m$,
ce qui, par définition de la dimension, signifie que
$\dim(A/\mathfrak p)=1$.

\begin{proposition}[10.5.10]
\label{IV.10.5.10-fr}
Soit $A$ un anneau local noethérien réduit et complet et qui n'est pas un
corps. Alors les intersections finies des noyaux des homomorphismes locaux de
$A$ dans des anneaux de valuation discrète $V$, qui font de $V$ une
$A$-algèbre finie, forment une base de filtre tendant vers $0$ pour la
topologie adique de $A$.
\end{proposition}

Il suffit (Bourbaki, \emph{Alg. comm.}, chap.~III, \S~2,
n\textsuperscript{o}~7, prop.~8) de prouver que l'intersection des noyaux
envisagés dans l'énoncé est réduite à $0$. Supposons d'abord que $A$ soit
intègre et de dimension $1$~; en vertu du théorème de Nagata ($\mathbf 0$,
23.1.5 et 23.1.6), la clôture intégrale $A'$ de $A$ est un anneau local
complet intègre et

